\chapter{比例、速度与平均数}\label{ch:g8:speed}

比例（\cref{ch:g7:prop}）现在走进了现实世界：
速度、流量、单位换算 --- 都是另外两个量的
\emph{商}。本章以加权平均数收尾，
其中比例思想可以避免一个经典错误。

\section{平均速度}

\begin{definition}[平均速度]\label{def:g8:speed:def}
某段行程的\emph{平均速度}\index{平均速度}是路程与时间的商，
\[
v = \frac{d}{t}
\qquad
\left(\text{路程除以时长}\right),
\]
以千米每小时（km/h）、米每秒（m/s）等为单位，\dots\ 
公式可以变形：$d = v \times t$ 与 $t = \frac{d}{v}$。
\end{definition}

\begin{example}\label{ex:g8:speed:basic}
一列火车行驶 $270$ km，用时 $1$ h $30$ min。先把时间化为小时：
$1$ h $30$ min $= 1.5$ h。然后
\[
v = \frac{270}{1.5} = 180 \text{ km/h}.
\]
以 $90$ km/h 的速度行驶 $2$ h $20$ min $= \frac73$ h 时的路程：
$d = 90 \times \frac73 = 210$ km。行驶 $35$ km、速度为 $14$ km/h 时所用的时间：
$t = \frac{35}{14} = 2.5$ h $= 2$ h $30$ min。
\end{example}

\begin{remark}[分钟不是小数]
$1$ h $30$ min 等于 $1.5$ h，而 $1$ h $20$ min 并\emph{不}等于 $1.2$ h：
它等于 $1 + \frac{20}{60} = \frac43 \approx 1.33$ h。
做除法前，一律先把分钟化为小时的分数（$60$ min $= 1$ h）。
\end{remark}

\begin{omfigure}
\begin{tikzpicture}
\begin{axis}[omaxis,
  width=8.4cm, height=5.6cm,
  xmin=0, xmax=3.4, ymin=0, ymax=270,
  xlabel={时间 (h)}, ylabel={路程 (km)},
  xtick={1,2,3}, ytick={80,160,240}]
  \addplot[omDef, thick, domain=0:3.2] {80*x};
  \addplot[omThm, thick, domain=0:3.2] {50*x};
  \node[omDef, font=\small, rotate=36] at (axis cs:2.35,215)
    {$80$ km/h};
  \node[omThm, font=\small, rotate=25] at (axis cs:2.6,155)
    {$50$ km/h};
\end{axis}
\end{tikzpicture}

{\small 速度恒定时，路程与时间成正比例：图像是一条
经过原点的直线，其陡峭程度\emph{即}为
速度。}
\end{omfigure}

\begin{method}[速度的换算]\label{met:g8:speed:convert}
将 km/h 化为 m/s：由 $1$ km $= 1000$ m 及 $1$ h $= 3600$ s，
\[
v \text{ km/h} = \frac{v \times 1000}{3600} \text{ m/s}
= \frac{v}{3.6} \text{ m/s}.
\]
将 m/s 化为 km/h 时，乘以 $3.6$。
\end{method}

\begin{example}\label{ex:g8:speed:convert}
$90$ km/h $= \frac{90}{3.6} = 25$ m/s。一名短跑运动员
跑 $100$ m 用时 $10$ s，速度为 $10$ m/s $= 36$ km/h。
\end{example}

\section{商的量}

\begin{example}[流量、密度、每千克价格]\label{ex:g8:speed:quotients}
速度只是同类量中的一种，它们的处理方法完全相同：
\begin{itemize}
  \item 一个水龙头注满 $48$ L 用时 $4$ min：流量
        $\frac{48}{4} = 12$ L/min，因此注满 $150$ L 的浴缸
        所需时间为 $\frac{150}{12} = 12.5$ min；
  \item $0.6$ kg 奶酪售价 $9$ 欧元：单价
        $\frac{9}{0.6} = 15$ 欧元每千克；
  \item 一辆汽车耗油 $6.3$ L，行驶 $90$ km：油耗
        $\frac{6.3}{90} \times 100 = 7$ L 每 $100$ km。
\end{itemize}
先找出其中的商，再运用三式关系
（$q = \frac ab$，$a = qb$，$b = \frac aq$）即可完成其余部分。
\end{example}

\section{加权平均数}

\begin{definition}[加权平均数]\label{def:g8:speed:average}
若数值 $v_1, v_2, \dots$ 各带有频数（或权重），
$n_1, n_2, \dots$，则其\emph{加权平均数}\index{加权平均数}
为
\[
\bar v = \frac{n_1 v_1 + n_2 v_2 + \dots}{n_1 + n_2 + \dots} :
\]
即各数值之和除以各权重之和。
\end{definition}

\begin{example}\label{ex:g8:speed:average}
一次测验由两组学生参加：第 1 组 $12$ 人，平均分
$11$；第 2 组 $18$ 人，平均分 $14$。全班的
平均分为：
\[
\bar v = \frac{12 \times 11 + 18 \times 14}{12 + 18}
= \frac{132 + 252}{30} = \frac{384}{30} = 12.8 .
\]
\emph{并非} $\frac{11 + 14}{2} = 12.5$：人数较多的一组会把
平均数拉向自己的数值——权重很重要。
\end{example}

\begin{example}[两段行程的平均速度]\label{ex:g8:speed:twolegs}
一名骑车人骑行 $30$ km，速度为 $30$ km/h，然后
骑行 $30$ km，速度为 $15$ km/h。全程的平均速度\emph{并非}
$\frac{30 + 15}{2} = 22.5$ km/h。应当按定义计算：
\begin{enumerate}
  \item 时间：$\frac{30}{30} = 1$ h，然后 $\frac{30}{15} = 2$ h；
        共 $3$ h；
  \item 总路程 $60$ km，因此
        $v = \frac{60}{3} = 20$ km/h。
\end{enumerate}
速度较慢的一段用时更长，因此权重更大 --- 速度的平均数
必须以\emph{时间}为权重。
\end{example}

\section{练习}

\begin{exercise}[$\star$]\label{exo:g8:speed:1}
求平均速度：$150$ km 用时 $2$ h；$27$ km 用时 $45$ min；
$100$ m 用时 $12.5$ s。
\end{exercise}

\begin{exercise}[$\star$]\label{exo:g8:speed:2}
求路程：行驶 $2$ h，速度为 $85$ km/h；行驶 $40$ min，速度为 $90$ km/h。
求时间：行驶 $315$ km，速度为 $90$ km/h（用 h 和 min 表示）。
\end{exercise}

\begin{exercise}[$\star$]\label{exo:g8:speed:3}
换算：$72$ km/h 化为 m/s；$5$ m/s 化为 km/h；$108$ km/h 化为 m/s。
\end{exercise}

\begin{exercise}[$\star$]\label{exo:g8:speed:4}
一个水龙头每分钟出水 $15$ L。注满一个 $600$ L 的水箱需要多长时间？
在 $2$ h $30$ 内能出水多少升？
\end{exercise}

\begin{exercise}[$\star$]\label{exo:g8:speed:5}
哪一种更划算：$1.2$ kg 苹果售价 $3.30$ 欧元，还是
$0.8$ kg 苹果售价 $2.16$ 欧元？请比较每千克的价格。
\end{exercise}

\begin{exercise}[$\star$]\label{exo:g8:speed:6}
一个班有 $25$ 名学生，一次测验的平均分为 $12.4$；另一
个班有 $35$ 名学生，平均分为 $10$。求这两个班合在一起的
平均分。
\end{exercise}

\begin{exercise}[$\star$]\label{exo:g8:speed:7}
声音约以 $340$ m/s 的速度传播。你看到闪电，
$6$ 秒后听到雷声。闪电击中的地方距你多远
（精确到 $100$ m）？
\end{exercise}

\begin{exercise}[$\star\star$]\label{exo:g8:speed:8}
一名徒步者步行 $2$ h，速度为 $5$ km/h，休息 $30$ min，然后步行 $1$ h $30$
，速度为 $4$ km/h。求总路程，以及
\emph{包括休息在内}的平均速度（总路程除以总经过时间）。
\end{exercise}

\begin{exercise}[$\star\star$]\label{exo:g8:speed:9}
某学生各科成绩及其系数：$15$（系数 $3$）、$9$
（系数 $2$）和 $12$（系数 $1$）。求加权
平均数。若不计权重，她的平均数是多少，以及
两者为何不同？
\end{exercise}

\begin{exercise}[$\star\star$]\label{exo:g8:speed:10}
一辆汽车行驶 $120$ km，速度为 $80$ km/h，然后行驶 $120$ km，速度为 $120$ km/h。
\begin{enumerate}
  \item 求每一段的用时，再求全程的
        平均速度。
  \item 解释为什么答案小于 $100$ km/h，
        即两个速度的中点。
\end{enumerate}
\end{exercise}

\begin{exercise}[$\star\star\star$]\label{exo:g8:speed:11}
两个村庄相距 $18$ km。安娜于 $10{:}00$ 从第一个村庄出发，
以 $5$ km/h 的速度步行前往第二个村庄；鲍里斯于
同一时刻从第二个村庄出发，以 $4$ km/h 的速度步行前往第一个村庄。他们
在什么时刻相遇，相遇处距安娜的村庄多远？（两人合起来
以 $5 + 4$ km/h 的速度缩小距离。）
\end{exercise}

\section{问题：调和平均数，或为何往返
会拉低平均速度}

\begin{problem}[{周末问题 --- 等距离上的平均速度
为 $\frac{2 v_1 v_2}{v_1 + v_2}$，且它从不超过两个速度的
中点}]
\label{pb:g8:speed:1}
\Cref{ex:g8:speed:twolegs} 与 \cref{exo:g8:speed:10} 都得出了同一个出人意料的结论：
以一个速度去，以另一个速度回，平均速度
\emph{低于}两者的中点。
本题要找出这一结论背后的精确公式，
--- 一种新的平均数，即
\emph{调和平均数}\index{调和平均数} --- 证明这一结论是一个定理，
并将其推到逻辑的极端：一场
世界上任何速度都无法实现的追赶。

\medskip
\noindent\textbf{第一部分 --- 往返的公式。}
一段行程中，距离 $d$ 以速度 $v_1$ 行走，然后
同样的距离 $d$ 以速度 $v_2$ 再行走（所有数均为正数）。
\begin{enumerate}
  \item 将两段用时以及总用时表示为
        含 $d$、$v_1$、$v_2$ 的分数。
  \item 全程的平均速度为
        $\bar v = \dfrac{2d}{\ \dfrac{d}{v_1} + \dfrac{d}{v_2}\ }$
        --- 一个繁分数
        （\cref{ex:g8:fractions:complex}）。化简它并证明：
        \[
        \bar v = \frac{2\, v_1 v_2}{v_1 + v_2} .
        \]
  \item 用上面引用的两个计算检验该公式：
        取 $v_1 = 30$，$v_2 = 15$ km/h
        （\cref{ex:g8:speed:twolegs}），并取 $v_1 = 80$，
        $v_2 = 120$ km/h（\cref{exo:g8:speed:10}）。
  \item 公式中的距离 $d$ 消失了。这对骑车人而言，
        具体意味着什么？
  \item 现在假设行程改为在每个速度下各花费相同的
        \emph{时间} $T$，证明此时平均速度为
        $\frac{v_1 + v_2}{2}$，即普通的中点。
        在\cref{def:g8:speed:average}的语境中：
        速度的平均数总要以什么量为权重，
        两种情形之间有何变化？
\end{enumerate}

\medskip
  \noindent\textbf{第二部分 --- 调和平均数对比算术平均数。}
设 $v_1$、$v_2$ 为两个正数，记
\[
H = \frac{2\, v_1 v_2}{v_1 + v_2}
\quad\text{（其\emph{调和平均数}），}
\qquad
A = \frac{v_1 + v_2}{2}
\quad\text{（其\emph{算术平均数}）。}
\]
\begin{enumerate}[resume]
  \item 计算 $H$ 和 $A$，数对分别取 $(30, 15)$ 和
        $(80, 120)$。每次都是哪一种平均数更大？
  \item 将 $A - H$ 通分，公分母取 $2(v_1 + v_2)$，
        用乘法公式展开分子
        （见\cref{pb:g8:equations:1}），并证明：
        \[
        A - H = \frac{(v_1 - v_2)^2}{2\,(v_1 + v_2)} .
        \]
  \item 由此推出所有这些“意外”背后的定理：对正速度而言，
        $H \leq A$ 恒成立，且等号当且仅当
        $v_1 = v_2$ 时成立。
        为什么分子中的平方项就能定出结论？
  \item 证明优美的恒等式 $H \times A = v_1 \times v_2$：
        两种平均数之积恰好还原为原来的积。
  \item 一辆汽车行驶 $60$ km，速度为 $40$ km/h，再行驶
        $60$ km，速度为 $60$ km/h。先用公式预测平均速度，
        再用时间与总路程的方法来
        加以验证。
\end{enumerate}

\medskip
\noindent\textbf{第三部分 --- 不可能的追赶。}
\begin{enumerate}[resume]
  \item 一名骑车人报名参加 $30$ km 的比赛，希望平均速度达到
        $30$ km/h。比赛全程最多可以占用多少总时间？
        她前 $15$ km 以 $15$ km/h 的速度骑行，
        这段时间预算中还剩多少时间？
  \item 证明她所期望的平均速度已经在字面上不可能实现：
        若她后半程以 $90$ km/h 的速度骑行，求其平均速度，
        并说明为何即使速度任意大，
        她也仍然达不到
        $30$ km/h。
  \item 她把目标降低为 $20$ km/h。
        后 $15$ km 应以多大速度骑行才能达成？
        请用调和平均数公式检验你的答案。
  \item 单独开冷水龙头，$20$ min 可注满浴缸；单独开热水龙头，
        $30$ min 可注满。两龙头同时开时，每分钟注满浴缸的几分之几？
        所需时间又是多少？
        将你的答案与 $H(20, 30)$ 比较 ---
        你注意到了什么？
  \item 一架飞机沿同一航线往返飞行：顺风时速度为 $900$ km/h，
        逆风时速度为 $700$ km/h。求往返全程的平均速度，
        并用一句话解释无论风怎样吹，
        为何风总会使往返行程\emph{变长}，
        即往返总时间增加。
\end{enumerate}
\end{problem}
